% ---------------------------------------------------------------- abstract
\begin{abstract}
Recognising a knot in an image means reading a planar diagram (where the strands
cross and which strand is on top) and then computing its topology. Image-to-label
classifiers skip the diagram: they cannot name a knot outside their training census
and cannot say when an image does not show enough to decide. We split the problem.
A dense network predicts diagram structure (crossing centres, the over-strand axis,
cable masks and strand centre-lines); a geometric parser turns the maps into a
diagram, and exact knot invariants identify it. Checks that do not trust the network
(an Euler count of faces, a four-arm test at every crossing, agreement across
test-time views and across two centre-line resolutions) turn likely errors into
abstentions. The formulation is backbone-agnostic: with a pair of standard U-Nets it
answers \UPCov\% of photo-like images on \bench{} (\UPUnseenCov\% on unseen
backgrounds) with no wrong answer and abstains on every silhouette, where a 250-class
image classifier is right on \ClfBenchAcc\% and a segmentation network alone can name
nothing. On a blind set of \BlindN{} drawings and renders from sources never used
during development, scored once, it is right on \UPBlindCorrect{} images and wrong on
\UPBlindWrong{}; the errors are drawings whose under-strand gaps are one or two pixels
wide. A U-Net with a transformer bottleneck, trained on the same data, beats the U-Net pair on a
second blind set at every one of three seeds (\BUUfVtwo{} vs \BUUnetVtwo{} correct of 312, on average) and
on real pictures, with about as many wrong answers. We release the benchmark, the blind protocol, an annotation tool for real cable
photographs and all code.
\end{abstract}
